\documentclass[9pt]{article}

\usepackage[margin=1.28cm]{geometry}
\usepackage{amsmath,amssymb,bm}
\usepackage{lmodern}
\usepackage{microtype}
\usepackage{enumitem}
\usepackage{setspace}

\setlength{\parindent}{0pt}
\setlength{\abovedisplayskip}{5pt}
\setlength{\belowdisplayskip}{5pt}
\setlength{\abovedisplayshortskip}{3pt}
\setlength{\belowdisplayshortskip}{3pt}
\setlength{\parskip}{2pt}

\title{\vspace{-1.0cm}\textbf{Planning and MPC Equations}}
\author{AC Track Control}
\date{}

\begin{document}
\maketitle
\vspace{-0.5cm}
\small
\setstretch{0.96}

\section{Reference-line optimization}

The lateral offset field $o_i$ is optimized over a closed racing line while
penalizing curvature, boundary violations, and non-smooth offsets:

\begin{align}
J_{\mathrm{line}}
&= t_{\mathrm{lap}}
+ w_{\kappa}\sum_i \kappa_i^2\,\Delta s_i \nonumber\\
&\quad
+ w_c\sum_i
\max\!\left(0,\,c_{\mathrm{target}}-c_i\right)^2\Delta s_i
+ w_s\sum_i \left\|\Delta^2 o_i\right\|_2^2,
\label{eq:line-cost}
\end{align}

where $t_{\mathrm{lap}}$ is the static lap-time proxy, $\kappa_i$ is the local
curvature, $c_i$ is the distance from the racing line to the nearest track
boundary, and $\Delta s_i$ is the segment length.

\section{Speed envelope}

The curvature-limited speed is

\begin{equation}
v_{\kappa,i}
=
\operatorname{clamp}\!\left(
\sqrt{\frac{a_{\mathrm{lat}}}{|\kappa_i|}},
\ v_{\min},
\ v_{\max}
\right).
\label{eq:curvature-speed}
\end{equation}

The forward acceleration and backward braking passes enforce

\begin{align}
v_i^2 &\leq v_{i-1}^2 + 2a_x^{\mathrm{accel}}\Delta s_{i-1},
\label{eq:forward-pass}\\
v_i^2 &\leq v_{i+1}^2 + 2a_x^{\mathrm{brake}}\Delta s_i.
\label{eq:backward-pass}
\end{align}

The high-speed planning response time is

\begin{equation}
\tau_{\mathrm{eff}}
=
\tau_0
+ \max\!\left(0,\,v-200\right)g_{\mathrm{hs}},
\qquad
v \text{ in km/h},
\label{eq:high-speed-response}
\end{equation}

and the braking-limited speed satisfies

\begin{equation}
v_{\mathrm{allowed}}^2
\leq
v_{\mathrm{ahead}}^2
+ 2a_{\mathrm{brake}}d_{\mathrm{eff}},
\qquad
d_{\mathrm{eff}}
=
d-v\tau_{\mathrm{eff}}.
\label{eq:braking-limit}
\end{equation}

\section{Lateral MPC}

The Frenet state and control input are defined as

\begin{equation}
x=
\begin{bmatrix}
e_y & e_\psi & v_x & v_y & r
\end{bmatrix}^{\!\top},
\qquad
u=
\begin{bmatrix}
\delta & a_x
\end{bmatrix}^{\!\top}.
\label{eq:lateral-state}
\end{equation}

The speed-scheduled continuous-time model is

\begin{align}
\dot e_y &= v_x e_\psi + v_y,
\label{eq:frenet-ey}\\
\dot e_\psi &= r,
\label{eq:frenet-epsi}\\
\dot v_x &= a_x,
\label{eq:frenet-vx}\\
\dot v_y &= a_{33}(v_x)v_y + a_{34}(v_x)r + b_{31}(v_x)\delta,
\label{eq:frenet-vy}\\
\dot r &= a_{43}(v_x)v_y + a_{44}(v_x)r + b_{41}(v_x)\delta.
\label{eq:frenet-r}
\end{align}

The identified coefficients are scheduled by speed and augmented by the
load-dependent yaw-response factor $\gamma(|a_y|)$:

\begin{align}
a_{33}(v_x) &= -\frac{c_{\mathrm{lat}}(v_x)}{v_x},
&
a_{34}(v_x) &= \frac{c_{\mathrm{yaw,lat}}(v_x)}{v_x}-v_x,
\label{eq:identified-a1}\\
a_{43}(v_x) &= \frac{c_{\mathrm{yaw,vel}}(v_x)}{v_x},
&
a_{44}(v_x) &= \frac{c_{\mathrm{yaw,damp}}(v_x)}{v_x}.
\label{eq:identified-a2}
\end{align}

With the steering actuator state, the model is extended by

\begin{equation}
\dot\delta
=
\frac{\delta_{\mathrm{cmd}}-\delta}{\tau_{\mathrm{steer}}}.
\label{eq:steer-actuator}
\end{equation}

For a horizon of $N$ stages, the LMPC cost is

\begin{align}
J_{\mathrm{lat}}
&=
\sum_{k=0}^{N-1}
\left[
\left(x_k-x_{\mathrm{ref},k}\right)^{\!\top}
Q_k
\left(x_k-x_{\mathrm{ref},k}\right)
+ u_k^{\!\top} R_k u_k
+ \Delta u_k^{\!\top} R_{d,k}\Delta u_k
\right]
\nonumber\\
&\quad+
\left(x_N-x_{\mathrm{ref},N}\right)^{\!\top}
Q_N
\left(x_N-x_{\mathrm{ref},N}\right),
\label{eq:lateral-cost}
\end{align}

where

\begin{equation}
Q_k=
\operatorname{diag}\!\left(
q_y m_y(v_k),\,
q_\psi m_\psi(v_k),\,
q_{v_x},\,
q_{v_y},\,
q_r m_r(v_k)
\right),
\label{eq:lateral-q}
\end{equation}

and $m_y$, $m_\psi$, and $m_r$ are speed-dependent weight maps.

\section{Longitudinal MPC}

The response lag is discretized by

\begin{equation}
\alpha=\frac{\Delta t}{\tau_{\mathrm{response}}},
\qquad 0\leq\alpha\leq 1.
\label{eq:longitudinal-alpha}
\end{equation}

The incremental acceleration model is

\begin{align}
a_{k+1} &= (1-\alpha)a_k+\alpha u_k,
\label{eq:longitudinal-accel}\\
v_{k+1} &= v_k+\Delta t\left(a_{k+1}-a_{\mathrm{coast}}(v_k)\right).
\label{eq:longitudinal-speed}
\end{align}

The longitudinal cost combines speed tracking, acceleration effort, and jerk:

\begin{equation}
J_{\mathrm{lon}}
=
\sum_{k=0}^{N-1}
\left[
q_v\left(v_k-v_{\mathrm{ref},k}\right)^2
+ q_a\left(a_k-a_{\mathrm{ref},k}\right)^2
+ r_a\Delta a_k^2
+ r_j\left(\Delta^2 a_k\right)^2
\right].
\label{eq:longitudinal-cost}
\end{equation}

The acceleration and jerk constraints are

\begin{align}
-a_{\mathrm{brake,max}}
&\leq a_k \leq a_{\mathrm{accel,max}},
\label{eq:longitudinal-accel-constraint}\\
-j_{\mathrm{brake,max}}\Delta t
&\leq \Delta a_k \leq
+j_{\mathrm{accel,max}}\Delta t.
\label{eq:longitudinal-jerk-constraint}
\end{align}

\section{Pedal mapping}

The acceleration request is converted to throttle by

\begin{equation}
T
=
\operatorname{clamp}\!\left(
\frac{a_{\mathrm{propulsion}}}
{a_{\mathrm{cap}}(v)},
0,
1
\right),
\label{eq:throttle-map}
\end{equation}

and to brake by interpolating the measured deceleration grid:

\begin{equation}
B
=
\mathcal{I}_{v,p}\!\left(a_{\mathrm{decel}}\right),
\qquad
p\in\{0,0.15,0.30,0.45,0.60,0.75,0.90,1.00\}.
\label{eq:brake-map}
\end{equation}

\section{Steady-state axis-response consistency relation}

The normalized-axis response can be projected onto the steady-state relation

\begin{equation}
\frac{u v}{r}
=
\frac{L}{k_{\mathrm{eff}}}
+ \frac{K_{\mathrm{eff}}}{k_{\mathrm{eff}}}v^2,
\label{eq:steering-calibration}
\end{equation}

where $L$ is the wheelbase. The effective scale $k_{\mathrm{eff}}$ and
derived understeer gradient $K_{\mathrm{eff}}$ are not independently
identifiable physical steering parameters unless the road-wheel angle is
measured separately. They summarize the same normalized-axis response already
identified by the dynamic model and are used here as a consistency check.

\section{Linear model and steering-actuator augmentation}

For the five-state Frenet model, the continuous-time state matrix is

\begin{equation}
A_c=
\begin{bmatrix}
0 & v_x & 0 & 1 & 0\\
0 & 0 & 0 & 0 & 1\\
0 & 0 & 0 & 0 & 0\\
0 & 0 & 0 & a_{33} & a_{34}\\
0 & 0 & 0 & a_{43} & a_{44}
\end{bmatrix},
\qquad
B_c=
\begin{bmatrix}
0 & 0\\
0 & 0\\
0 & 1\\
b_{31} & 0\\
b_{41} & 0
\end{bmatrix}.
\label{eq:continuous-matrices}
\end{equation}

The speed-scheduled coefficients are

\begin{align}
a_{33}&=-\frac{c_{\mathrm{lat}}}{v_x},&
a_{34}&=\frac{c_{\mathrm{yaw,lat}}}{v_x}-v_x,\\
a_{43}&=\frac{c_{\mathrm{yaw,vel}}}{v_x},&
a_{44}&=\frac{c_{\mathrm{yaw,damp}}}{v_x}.
\end{align}

With an actuator state $\delta$, the augmented continuous model is

\begin{equation}
\begin{bmatrix}
\dot e_y\\
\dot e_\psi\\
\dot v_x\\
\dot v_y\\
\dot r\\
\dot\delta
\end{bmatrix}
=
A_{\mathrm{aug}}
\begin{bmatrix}
e_y\\
e_\psi\\
v_x\\
v_y\\
r\\
\delta
\end{bmatrix}
+
B_{\mathrm{aug}}
\begin{bmatrix}
\delta_{\mathrm{cmd}}\\
a_x
\end{bmatrix},
\end{equation}

where

\begin{equation}
A_{\mathrm{aug}}=
\begin{bmatrix}
0 & v_x & 0 & 1 & 0 & 0\\
0 & 0 & 0 & 0 & 1 & 0\\
0 & 0 & 0 & 0 & 0 & 0\\
0 & 0 & 0 & a_{33} & a_{34} & b_{31}\\
0 & 0 & 0 & a_{43} & a_{44} & b_{41}\\
0 & 0 & 0 & 0 & 0 & -1/\tau_{\mathrm{steer}}
\end{bmatrix},
\end{equation}

\begin{equation}
B_{\mathrm{aug}}=
\begin{bmatrix}
0 & 0\\
0 & 0\\
0 & 1\\
0 & 0\\
0 & 0\\
1/\tau_{\mathrm{steer}} & 0
\end{bmatrix}.
\end{equation}

The sampled model used by the QP is

\begin{equation}
A_d=I+A_{\mathrm{aug}}\Delta t,
\qquad
B_d=B_{\mathrm{aug}}\Delta t,
\end{equation}

with the discrete spectral radius bounded before QP assembly.

\end{document}
